% Suggested insertion point:
% Chapter 4, after the general execution pipeline description or at the end of Section 4.4.

\subsection{Real-Robot Perception-to-Execution Pipeline}
\label{subsec:real_robot_pipeline}

In addition to the simulated benchmark, the system was deployed on a physical tabletop
assembly setup to validate that the perception, planning, and execution modules can be
connected in a real robot loop. The real-robot setup consists of a UR5 manipulator equipped
with a Robotiq parallel-jaw gripper, an RGB-D camera for object observation, a picking area
containing separated LEGO Duplo bricks, and an assembly area in which the target structure is
constructed. The target product is specified by the same structured YAML representation used
in simulation. Each block entry defines the part type, color, target position, and target yaw in
the assembly coordinate frame. The planner converts this product-level specification into an
ordered list of pick-and-place tasks, which is then consumed by the perception and execution
modules.

The real-robot system follows the same high-level structure as the simulated pipeline, but the
implementation uses a calibrated robot-camera transformation and taught reference poses to
reduce the sensitivity of open-loop Cartesian execution. At the beginning of a run, the robot
moves to a fixed observation pose and captures an RGB-D image of the picking area. For each
target part, GroundingDINO proposes open-vocabulary detections, SAM produces instance
masks, and FoundationPose estimates the CAD-based 6D object pose. The estimated pose is
transformed into the robot base frame and converted into a small Cartesian offset relative to
the taught pre-grasp pose. Placement targets are generated from the planned product YAML
and expressed as offsets relative to the taught pre-place pose. This design keeps the physical
execution close to repeatedly tested taught configurations while still using visual pose
estimation for object-specific grasp offsets and yaw correction.

\paragraph{Asynchronous perception during execution.}
To reduce idle time between consecutive assembly steps, the real-robot implementation uses
an asynchronous perception-execution schedule. In the synchronous version, the robot waits at
the observation pose, the vision module estimates the next target, the robot executes the
corresponding pick-and-place action, and the same sequence is repeated for the following
part. This creates avoidable waiting time because the vision pipeline is inactive while the
robot is moving.

In the asynchronous version, the first RGB-D observation is captured once at the fixed
observation pose and then reused as a frozen observation for subsequent target queries. While
the robot executes the current pick-and-place action, the vision process estimates the 6D pose
of the next target from the frozen observation. The selected instance mask is stored after each
successful target selection and excluded from later detections, which prevents repeated
selection of the same physical brick when multiple identical bricks are present in the initial
image. This is particularly important for products containing several identical white 2x4
bricks. The next execution step can therefore start as soon as the current action finishes and
the corresponding vision output is available, reducing the sequential wait time between
assembly steps.

This optimization does not change the geometric task plan. It only changes the scheduling
between perception and execution. The approach is suitable for the current real-robot
validation because all loose parts are visible in the initial observation and the picking area is
treated as a known initial scene. For more dynamic settings, the same interface could be
extended with repeated observations or closed-loop re-detection after each manipulation
step.
